\section{Kết quả đạt được, hướng phát triển}

\subsection{Kỹ năng và kiến thức thu thập được}

\subsubsection{Những việc đã hoàn thành}

Trong suốt thời gian thực tập, em đã tiếp nhận và hoàn thành hơn 170 phiếu công việc (bao gồm các yêu cầu tính năng nhỏ, tác vụ cải tiến và sửa lỗi) tương ứng với các đầu việc kỹ thuật trên các kho mã nguồn của sản phẩm. Các hạng mục đóng góp chính bao gồm:

\begin{itemize}[leftmargin=1.2cm, itemsep=1pt, topsep=2pt]
\item Tham gia phát triển tính năng khóa phòng bảo trì: từ mô hình dữ liệu, quy tắc nghiệp vụ, API phía máy chủ đến giao diện người dùng, bản dịch đa ngôn ngữ và viết kiểm thử tự động, đồng thời cập nhật công thức tính công suất phòng để loại trừ phòng đang ngừng kinh doanh.
\item Nghiên cứu và hiện thực hai chỉ tiêu thống kê mới về lượt khách nghỉ qua đêm và tổng số đêm nghỉ, đóng gói phần tính toán thành lớp xử lý thuần túy với các trường hợp kiểm thử tự động bao phủ các tình huống biên.
\item Tái cấu trúc cơ chế phân quyền hủy đơn thành hai quyền độc lập, chuẩn bị tập lệnh chuyển đổi dữ liệu an toàn cho các cơ sở đang hoạt động.
\item Các hạng mục tối ưu hiệu năng: giảm số lượt truy vấn và hạ độ phức tạp tính toán khi xuất tệp sổ quỹ, xử lý hiện tượng nghẽn tài nguyên cơ sở dữ liệu khi xuất hóa đơn chi tiết, bổ sung thời gian sống (TTL) cho khóa phân tán Redis và áp dụng quy ước tránh tính toán lặp trên giao diện.
\item Xây dựng các phân hệ giao diện cho bản quản lý mới: trang tổng quan (dashboard) với các biểu đồ doanh thu, danh mục phòng, tài liệu hướng dẫn tích hợp kênh bán phòng trực tuyến và bảng tổng kết bàn giao ca.
\item Địa phương hóa và đa ngôn ngữ: bổ sung tiếng Hàn, tiếng Trung, tiếng Lào cho các màn hình, xử lý đồng bộ định dạng ngày tháng, đơn vị tiền tệ và múi giờ, cùng các bài kiểm thử tự động trên cả hai tầng hệ thống.
\end{itemize}

\subsubsection{Kiến thức kỹ thuật}

Về cơ sở dữ liệu, em học được phương pháp phân tích câu truy vấn thực tế thay vì phụ thuộc vào mã sinh tự động từ ORM, nhận diện các điểm nghẽn phổ biến như phép nối bảng dư thừa, lọc trùng trên tập dữ liệu lớn hay nạp trước quan hệ không cần thiết. Em cũng tích lũy kinh nghiệm viết thủ tục lưu trữ tối ưu và xây dựng các kịch bản chuyển đổi dữ liệu an toàn cho môi trường sản xuất.

Về tối ưu hiệu năng, em hiểu rằng phần lớn giải pháp trong thực tế không đòi hỏi thuật toán phức tạp mà bắt nguồn từ những nguyên tắc rất cơ bản: hạn chế thực hiện lặp lại các phép tính tĩnh, không vận chuyển dữ liệu dư thừa qua mạng và thay thế quét tuần tự bằng tra cứu theo bảng băm.

Về kiến trúc phần mềm, em hiểu được giá trị của việc phân tầng và nguyên tắc phân định quyền sở hữu dữ liệu: trong hệ thống mà em tham gia, nhóm ưu tiên đặt logic tính toán nghiệp vụ tại dịch vụ sở hữu dữ liệu nguồn, tránh để cùng một chỉ tiêu bị tính toán độc lập ở nhiều nơi. Về kiểm thử tự động, nhận thức của em chuyển biến từ việc xem đây là thủ tục sang nhìn nhận kiểm thử như một bản đặc tả hành vi hệ thống và là cơ chế bảo vệ giúp hạn chế lỗi hồi quy khi mã nguồn thay đổi. Về đa ngôn ngữ, em nhận thức sâu sắc rằng địa phương hóa không chỉ là dịch chuỗi giao diện, mà bao gồm chuẩn hóa định dạng thời gian, số học, tiền tệ và đặc biệt là ranh giới múi giờ địa phương.

\subsubsection{Kỹ năng làm việc}

Bên cạnh kiến thức chuyên môn, kỳ thực tập đem lại cho em những kỹ năng làm việc quan trọng trong môi trường doanh nghiệp:

Kỹ năng đặt câu hỏi kỹ thuật: Thay vì mất nhiều thời gian tự tìm kiếm trong bế tắc, em hình thành phương pháp làm việc chủ động: sau một khoảng thời gian tự nghiên cứu và thử nghiệm hợp lý, nếu chưa tìm ra hướng đi thì chủ động trao đổi với người hướng dẫn kèm theo các phương án đã thử và giả thuyết cụ thể.

Kỹ năng viết nội dung commit: Từ thói quen viết nội dung commit chung chung, em dần hình thành cách mô tả rõ mã phiếu công việc và nội dung thay đổi, giúp tiết kiệm thời gian cho quá trình rà soát và truy vết mã nguồn sau này.

Phối hợp cùng bộ phận kiểm tra chất lượng (QA): Em hiểu rằng việc nhận phản hồi lỗi từ môi trường kiểm thử là một phần tất yếu của quy trình phát triển. Bộ phận QA tiếp cận hệ thống từ góc nhìn khách quan của người dùng thật, giúp phát hiện các tình huống biên mà lập trình viên dễ bỏ sót.

Sẵn sàng điều chỉnh giải pháp chưa tối ưu: Qua lần phải bỏ phần mã đã làm để chuyển logic sang đúng phân hệ, em hiểu rằng cần ưu tiên tính đúng đắn của kiến trúc hệ thống thay vì cố giữ lại giải pháp chỉ vì đã dành nhiều thời gian xây dựng.

\subsubsection{Những bài học em cho là quan trọng nhất}

\begin{enumerate}[leftmargin=1.2cm, itemsep=2pt, topsep=2pt]
\item Một tính năng chỉ thực sự hoàn thành khi nó hoạt động chuẩn xác và không gây ảnh hưởng tiêu cực đến các hành vi nghiệp vụ hiện có trên hệ thống.
\item Phạm vi của một thay đổi quan trọng tương đương nội dung của nó: với một lỗi cục bộ, cần ưu tiên giải pháp sửa đổi an toàn, có phạm vi hẹp thay vì vội vàng tái cấu trúc diện rộng khi chưa có kiểm thử bảo vệ.
\item Sử dụng ORM giúp tăng tốc độ phát triển nhưng không có nghĩa là có thể bỏ qua câu lệnh SQL thực tế được sinh ra; ở các tác vụ dữ liệu lớn, việc xem xét kế hoạch thực thi là cần thiết để phát hiện và xử lý các điểm nghẽn.
\item Quyền thao tác của người dùng phải được kiểm tra ở phía máy chủ, thay vì chỉ ẩn hoặc vô hiệu hóa nút trên giao diện; các ràng buộc trên giao diện chỉ hỗ trợ trải nghiệm người dùng chứ không thể thay thế cho việc kiểm soát quyền thực tế.
\item Kiểm thử tự động có giá trị lớn nhất trong việc bảo vệ hành vi hiện hữu của hệ thống; các ca kiểm thử giúp em tự tin hơn khi sửa đổi mã nguồn và hạn chế tối đa lỗi hồi quy phát sinh.
\end{enumerate}

\subsubsection{Những điều em chưa làm được}

Nhìn nhận một cách khách quan, em nhận thấy bản thân còn một số hạn chế cần tiếp tục hoàn thiện. Năng lực ước lượng thời gian thực hiện công việc đôi khi chưa sát với thực tế, chủ yếu do chưa lường hết thời gian cần thiết để đọc hiểu mã nguồn kế thừa. Việc viết kiểm thử tự động chưa thực sự đồng đều giữa các tầng: em tập trung nhiều vào logic nghiệp vụ phía máy chủ nhưng còn hạn chế ở tầng hiển thị giao diện. Em vẫn cần rèn luyện thêm khả năng tự chất vấn đa chiều các trường hợp ngoại lệ thay vì trông đợi vào sự định hướng của người hướng dẫn. Cuối cùng, mức độ hiểu biết của em về hệ thống mới tập trung sâu vào các phân hệ trực tiếp đảm nhiệm chứ chưa bao quát toàn diện mọi nhánh dịch vụ của sản phẩm.

\subsection{Hướng phát triển tiếp theo}

Đối với tính năng khóa phòng bảo trì, hướng phát triển tiếp theo là mở rộng cơ chế khóa phòng lặp lại theo chu kỳ định kỳ tự động và liên kết với kế hoạch quản lý tài sản, giúp chủ cơ sở theo dõi toàn diện lịch sử sửa chữa của từng phòng; mô hình sinh lệnh theo trạng thái thay đổi đã xây dựng có thể tái sử dụng trực tiếp cho yêu cầu này.

Đối với phân hệ thống kê lưu trú, giải pháp nâng cao hiệu năng trong dài hạn là áp dụng cơ chế tổng hợp trước (pre-aggregation) theo từng ngày và lưu vào bảng tổng hợp đệm; khi người dùng tra cứu khoảng thời gian tùy chọn, hệ thống chỉ cần cộng dồn các bản ghi ngày đã xử lý mà không phải quét lại toàn bộ dữ liệu gốc. Với công tác tối ưu hiệu năng, hướng đi tiếp theo là thiết lập hệ thống giám sát hiệu năng ứng dụng (Application Performance Monitoring - APM) để phát hiện sớm các truy vấn chậm trước khi phát sinh phản ánh từ khách hàng.

Về định hướng cá nhân, em đặt mục tiêu tiếp tục phát triển theo hướng kỹ sư toàn diện có khả năng bao quát cả tầng giao diện lẫn tầng máy chủ. Ba mục tiêu cụ thể em tập trung theo đuổi gồm: nghiên cứu chuyên sâu về tối ưu chỉ mục và kế hoạch thực thi câu truy vấn cơ sở dữ liệu; rèn luyện phương pháp phát triển hướng kiểm thử (Test-Driven Development - TDD) cho các phân hệ nghiệp vụ phức tạp; và nâng cao năng lực phân tích tác động toàn diện trước khi bắt tay vào hiện thực mã nguồn.

Kỳ thực tập tại doanh nghiệp không chỉ củng cố kiến thức học đường, mà quan trọng hơn đã trang bị cho em phương pháp làm việc chuẩn mực trên một hệ thống thương mại thực tế đang vận hành và phục vụ hàng nghìn người dùng mỗi ngày.
